% =====================================================================
%  Stage 1, раздел 2.
%  Сравнительный анализ публично доступных датасетов NIDS
%  и обоснование выбора UNSW-NB15.
% =====================================================================

\section{Сравнительный анализ публично доступных датасетов NIDS}

\subsection{Критерии оценки качества датасетов}

В литературе по сетевой безопасности систематизированы критерии, по которым
оценивается пригодность датасета для задач обучения и тестирования
NIDS~\cite{gharib2016evalframework, sharafaldin2018toward, ferrag2020dlcyber}.
В настоящей работе принят следующий набор критериев, объединяющий упомянутые
источники.

\begin{enumerate}[label=K\arabic*., leftmargin=2.5em]
    \item \textbf{Актуальность.} Соответствие представленных типов атак
    современной модели угроз: наличие классов \emph{Exploits, Backdoors,
    Reconnaissance, Botnet, DDoS}; временной горизонт записи трафика;
    репрезентативность прикладных сервисов.
    \item \textbf{Репрезентативность нормального трафика.} Наличие
    реалистичного фонового поведения с разнородными прикладными протоколами.
    \item \textbf{Качество разметки.} Соответствие меток реальной природе
    трафика; согласованность таксономии; отсутствие систематических ошибок
    маркировки.
    \item \textbf{Дисбаланс классов.} Доля атакующих наблюдений и её
    распределение по типам атак; чрезмерный или искусственно сбалансированный
    дисбаланс одинаково ухудшает обобщающую способность.
    \item \textbf{Пригодность для DL.} Объём, плотность временных меток,
    наличие признаковых представлений, допускающих формирование скользящих
    окон и последовательностей.
    \item \textbf{Пригодность для flow-анализа.} Наличие готовых flow-признаков
    или возможность их извлечения; согласованность со стандартами NetFlow/IPFIX.
    \item \textbf{Воспроизводимость.} Открытость доступа, фиксированность
    официальных train/test-разделений, наличие публикаций с описанием схемы
    сбора.
    \item \textbf{Ограничения и известные дефекты.} Зафиксированные в
    литературе аномалии конкретного датасета, способные исказить результаты
    обучения и оценки.
\end{enumerate}

В дальнейшем сравнении каждый датасет описывается с точки зрения этих
критериев, после чего проводится сводная оценка.

\subsection{KDDCup99}

KDDCup99~\cite{tavallaee2009nslkdd} --- историческое расширение трафика
DARPA 1998. Содержит 41 признак (включая host-based и time-based агрегации)
и четыре крупных класса атак (\emph{DoS, Probe, R2L, U2R}).

\begin{description}[font=\normalfont\itshape, leftmargin=1.5em]
    \item[Сильные стороны:] длительная история использования; сравнимость с
    большим числом ранних работ; малый объём, удобный для отладки.
    \item[Ограничения:] устаревшая модель угроз (атаки 1998 г.); существенное
    дублирование записей в train- и test-наборах; искусственная природа
    нормального трафика, не соответствующая корпоративным сценариям;
    статистические искажения, документированные в~\cite{tavallaee2009nslkdd,
    sommer2010outside}.
\end{description}

Согласно установившейся в литературе оценке~\cite{moustafa2016eval,
ferrag2020dlcyber}, использование KDDCup99 приводит к завышенным метрикам
(близким к насыщению) и не отражает современных классов атак, что делает
датасет малопригодным для обучения промышленно ориентированной системы.

\subsection{NSL-KDD}

NSL-KDD~\cite{tavallaee2009nslkdd} --- очищенная версия KDDCup99, полученная
устранением дубликатов и балансировкой долей классов в train/test-наборах.

\begin{description}[font=\normalfont\itshape, leftmargin=1.5em]
    \item[Сильные стороны:] устранены наиболее очевидные искажения KDDCup99;
    компактный объём; удобство для классических ML-методов и для
    воспроизводимых учебных экспериментов.
    \item[Ограничения:] таксономия и природа трафика идентичны KDDCup99 ---
    атаки и сервисы 1998 г.; не содержит современных классов атак (Exploits,
    Backdoors, Botnet); ограниченная пригодность для глубинных моделей из-за
    относительно небольшого объёма (\(\approx 125\) тыс. записей в KDDTrain+).
\end{description}

NSL-KDD принят в качестве вспомогательного датасета только для исторического
сравнения с литературой; в качестве основного обучающего корпуса он не
рассматривается.

\subsection{UNSW-NB15}

UNSW-NB15~\cite{moustafa2015unsw, moustafa2016eval, unswportal} был
сформирован на стенде \emph{IXIA PerfectStorm} в Cyber Range Lab Australian
Centre for Cyber Security. Содержит \(\approx 2{,}54\,\text{млн}\) записей
сетевых соединений, описанных 49 признаками, и девять семейств атакующего
трафика: \emph{Fuzzers, Analysis, Backdoors, DoS, Exploits, Generic,
Reconnaissance, Shellcode, Worms}~\cite{moustafa2015unsw}.

\begin{description}[font=\normalfont\itshape, leftmargin=1.5em]
    \item[Сильные стороны:]
    \begin{itemize}
        \item Современная (на момент публикации) таксономия атак, включающая
        классы \emph{Exploits} и \emph{Backdoors}, которых нет в KDD-семействе.
        \item Реалистичный нормальный трафик, генерируемый в условиях
        полнопроточного эмулятора прикладных сервисов.
        \item Признаки рассчитаны как \emph{flow-features} (длительность,
        счётчики байтов и пакетов, флаги TCP, межпакетные интервалы),
        совместимые со стандартом IPFIX~\cite{rfc7011}.
        \item Зафиксированы официальные разделения \emph{UNSW\_NB15\_training-set}
        и \emph{UNSW\_NB15\_testing-set}, упрощающие сравнение между
        работами~\cite{moustafa2016eval}.
        \item Богатый набор связанных публикаций
        автора~\cite{moustafa2016eval, koroniotis2019botiot,
        moustafa2021toniot}, обеспечивающий преемственность таксономии.
    \end{itemize}
    \item[Ограничения:]
    \begin{itemize}
        \item Естественный дисбаланс классов с доминированием класса
        \emph{Generic} в составе атакующих записей.
        \item Часть признаков (\emph{ct\_*}) рассчитывается как агрегаты по
        большому контекстному окну, что ограничивает их интерпретируемость в
        реальном времени.
        \item Часть мультиклассовых меток требует дополнительной валидации;
        фигурируют единичные противоречия в распределении \emph{Analysis} ---
        это отмечено в обзорных
        публикациях~\cite{moustafa2016eval}.
    \end{itemize}
\end{description}

UNSW-NB15 удовлетворяет всем восьми критериям K1--K8 в наибольшей степени и
выбран в работе как основной обучающий и тестирующий корпус
(\secref{subsec:unsw-justification}).

\subsection{CICIDS2017}

CICIDS2017~\cite{sharafaldin2018toward} был сформирован в Canadian Institute
for Cybersecurity (CIC) на стенде с эмуляцией поведения 25 пользователей и
содержит трафик за пять рабочих дней. Атаки представлены классами
\emph{Brute Force, Heartbleed, Botnet, DoS, DDoS, Web Attack, Infiltration,
Port Scan}.

\begin{description}[font=\normalfont\itshape, leftmargin=1.5em]
    \item[Сильные стороны:] современная таксономия с акцентом на прикладные
    атаки (Brute Force, Web Attack); подробная разметка по дням; flow-признаки
    извлекаются официальной утилитой CICFlowMeter и совместимы с IPFIX.
    \item[Ограничения:] в литературе документированы дефекты разметки и
    дубликаты в исходных
    CSV-файлах~\cite{gharib2016evalframework, ferrag2020dlcyber};
    значительная доля атакующего трафика для отдельных классов.
\end{description}

CICIDS2017 в настоящей работе используется в режиме cross-evaluation для
проверки переносимости методики (см. \secref{subsec:unsw-vs-cic}).

\subsection{CSE-CIC-IDS2018}

CSE-CIC-IDS2018~\cite{sharafaldin2018toward} --- расширенная версия
CICIDS2017, размещённая в облачной инфраструктуре AWS, с увеличенным числом
пользователей-эмуляторов (50) и расширенным набором атак, включая botnet и
infiltration. Объём датасета на порядок больше CICIDS2017.

\begin{description}[font=\normalfont\itshape, leftmargin=1.5em]
    \item[Сильные стороны:] высокий объём данных, удобный для DL;
    более актуальные классы атак.
    \item[Ограничения:] унаследованы документированные дефекты CIC-семейства;
    высокая стоимость воспроизведения экспериментов из-за объёма.
\end{description}

\subsection{TON\_IoT}

TON\_IoT~\cite{moustafa2021toniot} ---  датасет той же исследовательской
группы UNSW Canberra Cyber, ориентированный на IoT и edge-инфраструктуры.
Содержит сетевой трафик, журналы Windows/Linux, телеметрию IoT-устройств.

\begin{description}[font=\normalfont\itshape, leftmargin=1.5em]
    \item[Сильные стороны:] современная (2019--2021) модель угроз;
    мультиисточниковый формат; преемственность с UNSW-NB15.
    \item[Ограничения:] узкая ориентация на IoT-сегмент, что ограничивает
    применимость к классической корпоративной LAN.
\end{description}

TON\_IoT в работе рассматривается как кандидат на расширение методики и
обсуждается в направлениях дальнейшего развития (раздел 4 главы 4 итоговой ВКР).

\subsection{Bot-IoT}

Bot-IoT~\cite{koroniotis2019botiot} --- датасет той же группы, акцентирующий
поведение ботнетов в IoT-сегменте. Содержит более 73 миллионов записей с
явным доминированием атакующего трафика.

\begin{description}[font=\normalfont\itshape, leftmargin=1.5em]
    \item[Сильные стороны:] большой объём; современная DDoS-таксономия.
    \item[Ограничения:] чрезмерная доля атакующего трафика (\(>97\%\)) делает
    датасет нерепрезентативным с точки зрения нормального профиля и требует
    специальной балансировки; ориентация на IoT.
\end{description}

\subsection{Kyoto 2006+}

Kyoto~\cite{song2011kyoto} --- датасет, собранный с продуктивных honeypot-узлов
сети Kyoto University. Содержит подмножество признаков, заимствованных из
KDDCup99, дополненное метками, основанными на реальных IDS-наблюдениях.

\begin{description}[font=\normalfont\itshape, leftmargin=1.5em]
    \item[Сильные стороны:] реальный (а не синтетический) трафик; длительный
    временной горизонт.
    \item[Ограничения:] узкое признаковое пространство, ограничивающее
    применимость DL-моделей; смещение в сторону сценариев honeypot, не
    отражающих корпоративную LAN.
\end{description}

\subsection{DARPA IDS Dataset}

DARPA IDS 1998/1999~\cite{lippmann2000darpa} --- исторический корпус,
послуживший основой для KDDCup99.

\begin{description}[font=\normalfont\itshape, leftmargin=1.5em]
    \item[Сильные стороны:] историческая значимость; первая публичная
    \emph{labeled}-выборка.
    \item[Ограничения:] устаревшая модель угроз (1998 г.); искусственный
    нормальный трафик; существенные методологические нарекания, изложенные
    в~\cite{sommer2010outside, moustafa2016eval}.
\end{description}

DARPA в качестве обучающего датасета сегодня неприемлем; в работе он
упоминается исключительно для обзорной полноты.

\subsection{LITNET-2020}

LITNET-2020~\cite{damasevicius2020litnet} --- датасет, собранный в реальной
магистральной сети Литвы (LITNET), содержащий аннотированные сетевые
flow-записи за 10 месяцев.

\begin{description}[font=\normalfont\itshape, leftmargin=1.5em]
    \item[Сильные стороны:] реальный продуктивный трафик; продолжительный
    временной горизонт.
    \item[Ограничения:] таксономия и объём отдельных классов несбалансированы;
    относительно слабое распространение в литературе; на момент написания
    работы --- ограниченное число публикаций с воспроизводимыми сравнениями.
\end{description}

\subsection{CICDDoS2019 (контекстная ссылка)}

CICDDoS2019~\cite{sharafaldin2019ddos} --- специализированный датасет,
посвящённый DDoS-атакам различных типов (UDP, SYN, MSSQL, NetBIOS, NTP, SNMP,
DNS-amp). Используется в работе как косвенная справочная таксономия при
формировании библиотеки атакующих сценариев для ИИ-злоумышленника.

\subsection{Сводное сравнение датасетов}

В \tabref{tab:datasets-compare} приведена сводная экспертная оценка
рассмотренных датасетов по критериям K1--K8. Оценки выставлены качественно по
шкале (\(-\), \(\circ\), \(+\)), где \(+\) обозначает соответствие критерию,
\(\circ\) --- частичное соответствие, \(-\) --- несоответствие. Источниками
оценок служат описанные выше публикации.

\begin{table}[H]
\centering
\caption{Сводное сравнение публично доступных датасетов NIDS}
\label{tab:datasets-compare}
\renewcommand{\arraystretch}{1.2}
\setlength{\tabcolsep}{4pt}
\small
\begin{tabularx}{\linewidth}{@{}lcccccccX@{}}
\toprule
\textbf{Датасет} & \textbf{K1} & \textbf{K2} & \textbf{K3} & \textbf{K4} &
\textbf{K5} & \textbf{K6} & \textbf{K7} & \textbf{K8: ключевые ограничения} \\
\midrule
KDDCup99~\cite{tavallaee2009nslkdd}            & $-$    & $-$ & $-$ & $\circ$ & $\circ$ & $-$    & $+$ & устаревшая модель угроз; дубликаты \\
NSL-KDD~\cite{tavallaee2009nslkdd}             & $-$    & $-$ & $\circ$ & $+$ & $\circ$ & $-$    & $+$ & таксономия 1998 г.; малый объём для DL \\
UNSW-NB15~\cite{moustafa2015unsw}              & $+$    & $+$ & $+$ & $\circ$ & $+$    & $+$    & $+$ & доминирование Generic; ct\_*-агрегаты \\
CICIDS2017~\cite{sharafaldin2018toward}        & $+$    & $+$ & $\circ$ & $\circ$ & $+$    & $+$    & $+$ & дефекты разметки; дубликаты \\
CSE-CIC-IDS2018~\cite{sharafaldin2018toward}   & $+$    & $+$ & $\circ$ & $\circ$ & $+$    & $+$    & $\circ$ & унаследованные дефекты; объём \\
TON\_IoT~\cite{moustafa2021toniot}             & $+$    & $\circ$ & $+$ & $\circ$ & $+$    & $+$    & $+$ & ориентация на IoT \\
Bot-IoT~\cite{koroniotis2019botiot}            & $\circ$& $-$ & $+$ & $-$    & $+$    & $+$    & $+$ & доля атак $>97\%$ \\
Kyoto 2006+~\cite{song2011kyoto}               & $\circ$& $\circ$ & $\circ$ & $\circ$ & $\circ$ & $\circ$ & $\circ$ & honeypot-смещение; узкие признаки \\
DARPA~\cite{lippmann2000darpa}                 & $-$    & $-$ & $\circ$ & $\circ$ & $-$    & $\circ$ & $+$ & 1998 г.; искусственный нормальный трафик \\
LITNET-2020~\cite{damasevicius2020litnet}      & $+$    & $+$ & $\circ$ & $\circ$ & $+$    & $+$    & $\circ$ & таксономический дисбаланс \\
\bottomrule
\end{tabularx}
\end{table}

Из приведённой таблицы видно, что наибольшее покрытие критериев у двух
датасетов: UNSW-NB15 и CICIDS2017. UNSW-NB15 превосходит CICIDS2017 по
качеству разметки и преемственности таксономии (последняя продолжена в
TON\_IoT и Bot-IoT того же научного коллектива), а также по более ранней
доступности и большей укоренённости в обзорной
литературе~\cite{khraisat2019survey, ahmad2021nids, ferrag2020dlcyber}.

\subsection{Обоснование выбора UNSW-NB15 в качестве основного датасета}
\label{subsec:unsw-justification}

С учётом сводного сравнения и характера задачи диссертационной работы UNSW-NB15
выбран в качестве основного датасета по следующим причинам.

\begin{enumerate}
    \item \textbf{Соответствие модели угроз.} Девять семейств атак UNSW-NB15
    покрывают рассмотренные в \tabref{tab:traffic-anomalies} типы
    несанкционированных действий (DoS, Reconnaissance, Exploits, Backdoors,
    Worms), а класс \emph{Generic} даёт фоновую массу атак, удобную для
    оценки FP-rate в условиях дисбаланса.
    \item \textbf{Совместимость с IPFIX/NetFlow.} Признаки UNSW-NB15
    структурно эквивалентны полям IPFIX-записей, что обеспечивает прямую
    переносимость методики на инфраструктуры реальных корпоративных сетей.
    \item \textbf{Стабильность разделений.} Использование официальных
    разделений \emph{training-set/testing-set} обеспечивает воспроизводимость
    результатов и сопоставимость с литературой.
    \item \textbf{Преемственность таксономии.} Близость UNSW-NB15 к TON\_IoT и
    Bot-IoT позволит в направлениях дальнейшего развития систематически
    расширять методику на edge- и IoT-сегменты без перестройки label-схемы.
    \item \textbf{Достаточный объём для DL.} Объём примерно \(2{,}54\,\text{млн}\)
    записей удовлетворяет требованию K5 и согласуется с минимальными порогами
    объёмов, рекомендуемыми для обучения CNN--LSTM-моделей в DL-NIDS-обзорах.
\end{enumerate}

\subsection{UNSW-NB15 vs CICIDS2017: схема cross-evaluation}
\label{subsec:unsw-vs-cic}

Для проверки переносимости методики (требование, вытекающее из ограничения
K8 и из обзора~\cite{pendlebury2019tesseract}) принят следующий протокол.

\begin{itemize}
    \item Модель обучается на UNSW-NB15 (training-set).
    \item Производится унификация подмножества признаков, общих для UNSW-NB15
    и CICIDS2017 (\emph{duration}, \emph{bytes/pkts in/out}, \emph{flags},
    \emph{protocol}, \emph{service}); список общих признаков
    фиксируется в главе~2 настоящей работы.
    \item Тест проводится на размеченной выборке CICIDS2017 без дообучения и
    с дообучением на 10\% выборки CICIDS2017 (контрольный режим transfer
    learning).
    \item Оценивается падение F1, PR-AUC и FP-rate. Падение интерпретируется
    как мера переносимости (\eng{generalization gap}).
\end{itemize}

\subsection{Угрозы валидности и пути их минимизации}

С учётом замечаний~\cite{sommer2010outside, pendlebury2019tesseract,
gharib2016evalframework} в работе предусмотрены следующие меры по снижению
угроз внутренней и внешней валидности экспериментов:

\begin{enumerate}
    \item Использование официального train/test-разделения UNSW-NB15 без
    дополнительного перемешивания между разделами.
    \item Контроль временного смещения (\emph{temporal split}) при
    формировании скользящих окон --- запрет утечек ``из будущего''.
    \item Прогон экспериментов с фиксированными seed для воспроизводимости
    (\(\geq 5\) seed на эксперимент).
    \item Кросс-проверка результатов на CICIDS2017 как независимом датасете.
    \item Учёт стоимости ложных срабатываний на единицу времени
    (FP-per-time), что компенсирует потенциально завышенную ROC-AUC при
    дисбалансе.
\end{enumerate}

\subsection{Промежуточные выводы по разделу}

\begin{enumerate}
    \item Анализ десяти публично доступных корпусов трафика по восьми
    критериям K1--K8 показал, что наибольшее покрытие критериев имеют
    UNSW-NB15 и CICIDS2017; остальные либо устарели по модели угроз
    (KDD-семейство, DARPA), либо ориентированы на IoT-сегмент
    (TON\_IoT, Bot-IoT), либо имеют известные методологические дефекты
    разметки.
    \item UNSW-NB15 превосходит CICIDS2017 по качеству разметки,
    преемственности таксономии и совместимости со стандартом IPFIX,
    что обосновывает его выбор в качестве основного обучающего и тестирующего
    корпуса.
    \item CICIDS2017 принят как вспомогательный корпус для cross-evaluation;
    разработана схема унификации общих признаков и протокол оценки переносимости
    методики.
    \item Зафиксирован набор мер по минимизации угроз валидности
    (фиксированные разделения, временной split, повторение экспериментов
    по seed, кросс-проверка), что соответствует современным требованиям к
    обучающим протоколам в области ML-NIDS.
\end{enumerate}
